\subsection{Dataset preprocessing and construction}

The promoter prediction datasets used in this work were assembled from multiple
public sources with heterogeneous file formats, including FASTA files,
gzip-compressed FASTA files, and line-oriented text files. To ensure a uniform
experimental representation, all sources were first converted into a common
FASTA naming scheme encoding the class, source study, organism type, dataset
identifier, number of records, and original sequence length. Promoter sequences
were assigned label $1$, whereas non-promoter sequences were assigned label $0$.

All model-ready sequences were normalized to a fixed length of 81 bp around the
transcription start site (TSS). The fixed window contains 60 bp upstream of the
TSS, the TSS nucleotide itself, and 20 bp downstream. Because the original
sequence windows differ across studies, the TSS coordinate used for cropping was
defined per source dataset. GraphPro sequences of length 81 used TSS position
60; CNNProm sequences used position 60 for 81 bp windows and position 200 for
251 bp windows; CDBProm sequences used position 60 for 100 bp windows and
position 500 for 1001 bp windows; and Bhandari et al. sequences used position
700 for 1000 bp windows.

After cropping, promoter and non-promoter FASTA files were paired by source
study, organism type, and dataset identifier. Each pair was converted into CSV
files containing two fields: \texttt{sequence}, the 81 bp DNA sequence, and
\texttt{label}, the binary class label. The processed intra-study datasets use a
filename convention that records the source study, organism type, dataset name,
positive count, negative count, and partition.

Finally, a strict DNA cleaning step was applied to produce the model-ready
processed dataset directory. This step removed only rows whose sequence contained
symbols outside the alphabet \{A, C, G, T\}; no original processed CSV file was
overwritten. Across 58 processed CSV files, 338 rows were removed from 12 files.
After this cleaning step, all 58 cleaned CSV files passed validation for 81 bp
sequence length, binary labels, and strict DNA alphabet compliance.
